\section{Step 43 Repeat a stream with explicit chunk history}\label{fir-s43}

Split a sixty-four-output input after output twenty-three. The first call requests twenty-four outputs; the second requests forty and starts forty-eight bytes after the original input base. Its first seven samples are the history needed for the next current sample. Keep coefficients unchanged and concatenate the returned output samples. No retained kernel state should be needed.

\subsection{Expected observation}
The concatenated result equals a complete call. The checks verify all sixty-five cuts, including an empty first or second call, in one session.

\Needspace{11\baselineskip}
\subsection{Check your understanding}
Why does the second input begin at original sample position twenty-four?

\begin{itemize}
\item \textbf{A} That starts the next window with its seven required history samples.
\item \textbf{B} Every output consumes eight new input samples.
\item \textbf{C} The next call discards history automatically.
\end{itemize}
